\section{군 준동형과 구조를 보존하는 함수}
\label{sec:homomorphisms}

군 사이의 함수 중에서 연산을 보존하는 함수가 중요하다. 군의 원소 이름이 달라도 연산 구조가 같을 수 있으며, 준동형은 이 구조를 비교하는 도구이다.

\begin{definition}[군 준동형]
군 $G,H$ 사이의 함수 $\varphi:G\to H$가 모든 $a,b\in G$에 대해
\[
  \varphi(ab)=\varphi(a)\varphi(b)
\]
를 만족하면 $\varphi$를 \emph{군 준동형}이라고 한다.
\end{definition}

정의에는 항등원과 역원을 보존한다는 조건이 따로 없지만, 곱을 보존하면 자동으로 따라온다.

\begin{proposition}
군 준동형 $\varphi:G\to H$에 대하여
\[
  \varphi(e_G)=e_H,
  \qquad
  \varphi(a^{-1})=\varphi(a)^{-1},
  \qquad
  \varphi(a^n)=\varphi(a)^n
\]
가 모든 $a\in G$, $n\in\Z$에 대해 성립한다.
\end{proposition}

\begin{proof}
\[
  \varphi(e_G)=\varphi(e_Ge_G)=\varphi(e_G)^2
\]
이므로 $H$에서 소거하여 $\varphi(e_G)=e_H$를 얻는다. 또한
\[
  e_H=\varphi(e_G)=\varphi(aa^{-1})=
  \varphi(a)\varphi(a^{-1})
\]
이므로 역원의 유일성에 의해 $\varphi(a^{-1})=\varphi(a)^{-1}$이다. 거듭제곱에 관한 식은 양의 지수에 대한 귀납법과 음의 지수 정의로 따른다.
\end{proof}

\begin{example}[정수에서 시작하는 모든 준동형]
군 $G$의 원소 $g$를 하나 고정하면
\[
  \varphi_g:\Z\to G,\qquad n\mapsto g^n
\]
은 준동형이다. 반대로 $\Z$에서 $G$로 가는 모든 준동형은 $1$의 상 $g=\varphi(1)$에 의해 이 꼴로 결정된다.
\end{example}

\begin{example}[행렬식]
\[
  \det:\GL_n(K)\to K^\times
\]
는 $\det(AB)=\det(A)\det(B)$이므로 준동형이다.
\end{example}

\begin{example}[지수함수]
\[
  \exp:(\R,+)\to(\R_{>0},\times)
\]
는 $\exp(x+y)=\exp(x)\exp(y)$이므로 준동형이며, 실제로는 동형이다. 역함수는 자연로그이다.
\end{example}

\begin{definition}[핵과 상]
준동형 $\varphi:G\to H$의 \emph{핵}과 \emph{상}을 각각
\[
  \ker\varphi=\{g\in G:\varphi(g)=e_H\},
  \qquad
  \im\varphi=\{\varphi(g):g\in G\}
\]
로 정의한다.
\end{definition}

핵은 $H$에서 항등원으로 ``사라지는'' 원소들의 집합이고, 상은 실제로 도달하는 부분이다.

\begin{proposition}
$\ker\varphi\le G$이고 $\im\varphi\le H$이다.
\end{proposition}

\begin{proof}
$a,b\in\ker\varphi$이면
\[
 \varphi(ab^{-1})=\varphi(a)\varphi(b)^{-1}=e_He_H^{-1}=e_H
\]
이므로 부분군 판정법에 의해 $\ker\varphi\le G$이다. 상에 대해서는 $x=\varphi(a)$, $y=\varphi(b)$일 때
\[
  xy^{-1}=\varphi(a)\varphi(b)^{-1}=\varphi(ab^{-1})\in\im\varphi
\]
이므로 역시 부분군이다.
\end{proof}

\begin{proposition}[단사성과 핵]
군 준동형 $\varphi:G\to H$가 단사일 필요충분조건은
\[
  \ker\varphi=\{e_G\}
\]
이다.
\end{proposition}

\begin{proof}
$\varphi$가 단사이면 $\varphi(g)=e_H=\varphi(e_G)$에서 $g=e_G$이므로 핵은 자명하다. 반대로 핵이 자명하고 $\varphi(a)=\varphi(b)$라 하자. 그러면
\[
  \varphi(ab^{-1})=\varphi(a)\varphi(b)^{-1}=e_H
\]
이므로 $ab^{-1}\in\ker\varphi=\{e_G\}$이다. 따라서 $a=b$이다.
\end{proof}

\begin{example}
행렬식 준동형의 핵은
\[
  \ker(\det)=\SL_n(K)
\]
이다. 따라서 $\SL_n(K)$는 단순히 우연히 닫힌 행렬 집합이 아니라, 자연스러운 준동형이 정보를 잃어버리는 정확한 부분이다.
\end{example}

\subsection{동형과 자동동형}

\begin{definition}
전단사 군 준동형을 \emph{군 동형}이라고 한다. $G$와 $H$ 사이에 동형이 존재하면
\[
  G\cong H
\]
라고 쓴다. $G$에서 자기 자신으로 가는 동형을 \emph{자동동형}이라고 하며, 자동동형 전체의 군을 $\Aut(G)$라고 쓴다.
\end{definition}

동형인 군은 원소의 이름만 다를 뿐 군론적으로 같은 구조를 가진다. 군의 위수, 아벨성, 원소의 위수, 부분군 구조 등은 동형 아래에서 보존된다.

\begin{example}[내부자동동형]
$a\in G$에 대해
\[
  c_a:G\to G,\qquad c_a(g)=aga^{-1}
\]
는 자동동형이다. 이를 $a$에 의한 \emph{켤레} 또는 \emph{내부자동동형}이라고 한다. 군이 아벨이면 모든 내부자동동형은 항등함수이다.
\end{example}

\subsection{케일리 정리}

모든 군은 어떤 순열군의 부분군으로 볼 수 있다. 이 결과는 군을 ``추상적 기호''가 아니라 실제 변환으로 해석할 수 있음을 보장한다.

\begin{theorem}[케일리 정리]
모든 군 $G$는 $\Sym(G)$의 부분군과 동형이다. 특히 $|G|=n<\infty$이면 $G$는 $S_n$의 어떤 부분군과 동형이다.
\end{theorem}

\begin{proofidea}
각 $g\in G$가 $G$ 위에서 왼쪽 곱셈
\[
  L_g(x)=gx
\]
을 수행한다고 보자. $L_g$는 역함수 $L_{g^{-1}}$를 가지므로 순열이다. $g\mapsto L_g$가 단사 준동형임을 확인하면 된다.
\end{proofidea}

\begin{proof}
각 $g\in G$에 대해 $L_g:G\to G$, $x\mapsto gx$를 정의한다. $L_{g^{-1}}$가 역함수이므로 $L_g\in\Sym(G)$이다. 또한
\[
  (L_g\circ L_h)(x)=g(hx)=(gh)x=L_{gh}(x)
\]
이므로
\[
  L:G\to\Sym(G),\qquad g\mapsto L_g
\]
는 준동형이다. $L_g=L_h$이면 항등원 $e$에 대입하여 $g=L_g(e)=L_h(e)=h$를 얻으므로 $L$은 단사이다. 따라서 $G\cong L(G)\le\Sym(G)$이다.
\end{proof}

\begin{checkpoint}
\begin{enumerate}[label=\arabic*.]
  \item $\varphi:\Z\to\Z/12\Z$, $n\mapsto\overline n$의 핵과 상을 구하라.
  \item $\varphi:\R^\times\to\{\pm1\}$, $x\mapsto x/|x|$의 핵을 구하라.
  \item 준동형이 생성집합의 상에 의해 완전히 결정되는 이유를 설명하라.
  \item 케일리 정리의 매장 $G\hookrightarrow\Sym(G)$에서 $g\ne h$이면 왜 $L_g\ne L_h$인가?
\end{enumerate}
\end{checkpoint}
